\documentclass[12pt]{article}
\usepackage[T1]{fontenc}
\usepackage[utf8]{inputenc}
\usepackage[brazil]{babel}
\usepackage[margin=1in]{geometry}
\setlength{\parindent}{0pt}
\begin{document}
\section*{Tema 36}
Processo: PEDILEF 0052862-57.2008.4.03.6301/SP\\
Situacao: Julgado (Súmula 77 da TNU)\\
Ramo: DIREITO PREVIDENCIÁRIO\\
Questao: Saber se é necessário exame das condições pessoais do segurado considerado capaz para o trabalho, nos pedidos de concessão de auxílio-doença.\\
Tese: Na concessão do auxílio-doença é dispensável o exame das condições pessoais do segurado quando não constatada a incapacidade laboral. Vide Súmula 77 da TNU.\\
Fonte: https://www.cjf.jus.br/phpdoc/virtus/
\end{document}
